\section*{अध्याय \ref{ch:g8:fronts-and-weather-maps} --- वायु राशियाँ, वाताग्र और मौसम मानचित्र}

\begin{solution}{exo:g8:fronts-and-weather-maps:1}
(क) समुद्री उष्णकटिबंधीय: गरम और नम। (ख) महाद्वीपीय ध्रुवीय: ठंडी और सूखी।
(ग) महाद्वीपीय उष्णकटिबंधीय: बहुत गरम और सूखी। (घ) समुद्री ध्रुवीय: ठंडी और
नम।
\end{solution}

\begin{solution}{exo:g8:fronts-and-weather-maps:2}
त्रिभुजों वाली नीली रेखा: \omterm{def:g8:fronts-and-weather-maps:front}{शीत वाताग्र}। अर्धवृत्तों वाली लाल रेखा: \omterm{def:g8:fronts-and-weather-maps:front}{उष्ण वाताग्र}। एक जैसे
दाब वाली जगहों को जोड़ती रेखा: \omterm{def:g8:fronts-and-weather-maps:isobar}{समदाब रेखा}। ``उच्च'': \omterm{def:g8:fronts-and-weather-maps:depression}{प्रतिचक्रवात} का केंद्र (उच्च दाब)।
``निम्न'': \omterm{def:g8:fronts-and-weather-maps:depression}{अवदाब} का केंद्र (निम्न
दाब)।
\end{solution}

\begin{solution}{exo:g8:fronts-and-weather-maps:3}
(क) सही। (ख) ग़लत: वे उससे पहले आते हैं, बहुत आगे, आकाश में ऊँचे। (ग) ग़लत:
पास-पास की \omterm{def:g8:fronts-and-weather-maps:isobar}{समदाब रेखाओं} का अर्थ है तेज़ पवन। (घ) सही।
\end{solution}

\begin{solution}{exo:g8:fronts-and-weather-maps:4}
(क) \qty{980}{hPa}। (ख) \qty{1034}{hPa}। (ग) \qty{4}{hPa}। (घ) $1034 -
980 = \qty{54}{hPa}$।
\end{solution}

\begin{solution}{exo:g8:fronts-and-weather-maps:5}
\omterm{def:g4:clouds-rain-snow:cloud-families}{पक्षाभ मेघ} की पतली धारियाँ; मध्य \omterm{def:g4:clouds-rain-snow:cloud-families}{स्तरी मेघ} का एक धूसर घूँघट; \omterm{def:g4:clouds-rain-snow:cloud}{बादल} नीचे और घने होते
जाते हैं; लगातार \omterm{def:g4:clouds-rain-snow:precipitation}{बारिश}; ज़्यादा गरम, ज़्यादा नम हवा।
\end{solution}

\begin{solution}{exo:g8:fronts-and-weather-maps:6}
\omterm{def:g8:fronts-and-weather-maps:front}{शीत वाताग्र} पर घनी ठंडी हवा एक \emph{फन्नी} की तरह गरम हवा के नीचे घुसती है: गरम हवा
एक \emph{खड़ी} ढलान पर ऊपर उठाई जाती है, तेज़ी से उठती है और ऊँचे \omterm{def:g4:clouds-rain-snow:cloud}{बादल} बनाती है। \omterm{def:g8:fronts-and-weather-maps:front}{उष्ण
वाताग्र} पर गरम हवा ठंडी हवा के ऊपर एक \emph{हल्की ढलान} पर धीरे-धीरे चढ़ती है: वह एक
लंबी दूरी में थोड़ा-थोड़ा करके उठती है, और उसके \omterm{def:g4:clouds-rain-snow:cloud}{बादल} परतों में फैल
जाते हैं।
\end{solution}

\begin{solution}{exo:g8:fronts-and-weather-maps:7}
(क) ``निम्न'' के पास: वहाँ \omterm{def:g8:fronts-and-weather-maps:isobar}{समदाब रेखाएँ} ज़्यादा पास-पास हैं, और तीर ज़्यादा लंबे हैं।
(ख) पूर्व से (तीर पश्चिम की ओर, थोड़ा दक्षिण-पश्चिम की ओर, इशारा करते हैं)।
(ग) पूर्व से (तीर पश्चिम की ओर इशारा करते हैं)। (घ)
वामावर्त।
\end{solution}

\begin{solution}{exo:g8:fronts-and-weather-maps:8}
(क) $1016 - 1003 = \qty{13}{hPa}$। (ख) एक \omterm{def:g8:fronts-and-weather-maps:depression}{अवदाब}: दाब गिरता ही जा रहा है।
(ग) \omterm{def:g4:clouds-rain-snow:cloud}{बादल}, हवा और \omterm{def:g4:clouds-rain-snow:precipitation}{बारिश}, शायद एक \omterm{def:g8:fronts-and-weather-maps:front}{वाताग्र} के साथ।
\end{solution}

\begin{solution}{exo:g8:fronts-and-weather-maps:9}
(क) cT, महाद्वीपीय उष्णकटिबंधीय। (ख) बहुत गरम और सूखी। (ग) \omterm{def:g2:water-on-earth:ocean}{समुद्र} के ऊपर वह जलवाष्प
लेती है: वह ज़्यादा नम हो जाती है, और पानी के पास की हवा \omterm{def:g2:water-on-earth:ocean}{समुद्र} से थोड़ी ठंडी हो जाती
है।
\end{solution}

\begin{solution}{exo:g8:fronts-and-weather-maps:10}
(क) $4 - 1 = 3$ दिन। (ख) $3 \div 3 = 1$ दिन हर दस वर्षों में। (ग) उसे
$5 - 1 = 4$ दिन और चाहिए: लगभग 40 वर्ष।
\end{solution}

\begin{solution}{exo:g8:fronts-and-weather-maps:11}
(क) लगभग \qty{1350}{km} (साढ़े चार पैमाना-पट्टियाँ)। (ख) \qty{4}{hPa},
\qty{80}{km} में: $4 \div 80 \times 100 = \qty{5}{hPa}$ प्रति \qty{100}{km}।
(ग) $4 \div 140 \times 100 \approx \qty{2.9}{hPa}$ प्रति \qty{100}{km}।
(घ) ``निम्न'' के पास, जहाँ दाब ज़्यादा तेज़ी से बदलता है; हाँ, ``निम्न'' के पास तीर
ज़्यादा लंबे हैं।
\end{solution}

\begin{solution}{exo:g8:fronts-and-weather-maps:12}
(क) \omterm{def:g8:fronts-and-weather-maps:front}{शीत वाताग्र}: $300 \div 50 = 6$ घंटे। \omterm{def:g8:fronts-and-weather-maps:front}{उष्ण वाताग्र} \qty{25}{km} प्रति घंटा की चाल
से चलता है: $200 \div 25 = 8$ घंटे। \omterm{def:g8:fronts-and-weather-maps:front}{शीत वाताग्र} 2 घंटे पहले पहुँचता है। (ख) पहला
शहर: मीनार जैसे ऊँचे \omterm{def:g4:clouds-rain-snow:cloud}{बादल}, भारी बौछारें, झोंके, फिर ठंडी हवा और साफ़ होता आकाश। दूसरा
शहर: नीचे आते परतदार \omterm{def:g4:clouds-rain-snow:cloud}{बादल} और एक लंबी, लगातार \omterm{def:g4:clouds-rain-snow:precipitation}{बारिश}, फिर ज़्यादा गरम, ज़्यादा नम हवा।
(ग) दूसरा शहर, जहाँ गरम हवा पहुँचती है।
\end{solution}

\begin{solution}{pb:g8:fronts-and-weather-maps:1}
\textbf{1.} गरम हवा में, दोनों \omterm{def:g8:fronts-and-weather-maps:front}{वाताग्रों} के बीच। \omterm{def:g8:fronts-and-weather-maps:front}{उष्ण वाताग्र} पहले ही गुज़र चुका
है।

\textbf{2.} दक्षिण-पश्चिम से।

\textbf{3.} पूर्व की ओर: \omterm{def:g8:fronts-and-weather-maps:front}{शीत वाताग्र} के त्रिभुज उस ओर लगे होते हैं जिस ओर वह चलता है,
और वे पूर्व की ओर, शहर की ओर, इशारा करते हैं।

\textbf{4.} मीनार जैसे ऊँचे \omterm{def:g4:clouds-rain-snow:cloud}{बादल}, भारी बौछारें, शायद गरज और \omterm{def:g4:clouds-rain-snow:precipitation}{ओले}, हवा के झोंके; फिर
ठंडी हवा, पश्चिम या उत्तर-पश्चिम से पवन और साफ़ होता
आकाश।

\textbf{5.} 06:00 बजे \qty{180}{km} पर और 09:00 बजे \qty{300}{km} पर।

\textbf{6.} $300 - 180 = \qty{120}{km}$।

\textbf{7.} $120 \div 3 = \qty{40}{km}$ प्रति घंटा।

\textbf{8.} शहर \qty{480}{km} पर है: $480 - 300 = \qty{180}{km}$।

\textbf{9.} $180 \div 40 = 4.5$ घंटे, यानी 4 घंटे और 30 मिनट।

\textbf{10.} 09:00 + 4~घंटे~30~मिनट = \textbf{13:30}।

\textbf{11.} नहीं: बौछारें लगभग 14:30 बजे, मैच शुरू होने से पहले, ख़त्म हो जानी
चाहिए।

\textbf{12.} $180 \div 30 = 6$ घंटे: 15:00 बजे, ठीक मैच शुरू होने के समय।
पूर्वानुमान: ``\omterm{def:g8:fronts-and-weather-maps:front}{शीत वाताग्र} लगभग \textbf{13:30} बजे शहर तक पहुँचना चाहिए, एक घंटे की
भारी बौछारों और झोंकों के साथ, फिर आकाश साफ़ होगा: मैच शायद खेला जा सकेगा। अगर \omterm{def:g8:fronts-and-weather-maps:front}{वाताग्र}
धीमा हो जाता है, तो वह 15:00 बजे पहुँच सकता है और मैच में देर करा सकता
है।''
\end{solution}
